\section{2024 ROBOCON2024-颗粒归仓}

作为华东理工大学“无贰”战队的副队长与机械组组长，设计了一台参赛机器人的整车结构。

\begin{figure}[H]
    \centering
    \includegraphics[width=0.95\textwidth]{robocon-2024-1.jpg}
\end{figure}

2024年的ROBOCON赛事的主题为“颗粒归仓”，要求参赛队设计两台能够在规定时间内完成抓取、投掷与搬运任务的机器人，其中一台（R1）要求能在一区与二区完成夹取放置秧苗、夹取并传递谷物的任务；另一台（R2）要求能在三区全自动地抓取特定的谷物并放置入筒仓之中。

作为华东理工大学“无贰”战队的副队长与机械组组长，我主要负责R1的整车机械结构设计工作，同时也参与了R2的结构设计工作。此外还负责战队的日常管理与运营工作，包括在招新期间组织多次机械组培训，提高团队整体的技术水平和工作效率。

主要贡献包括：
\begin{itemize}
    \item \textbf{底盘结构设计：}设计了两套使用DJI M3508电机驱动的麦轮底盘，其中一套使用成品麦轮并通过M3508的电机轴输出，结构简单但对电机轴系有较大的弯矩载荷；另一套则使用内嵌结构设计的定制麦轮以提升底盘的紧凑性与轴系的稳定承载能力。
    \item \textbf{夹取与传递机构：}在R1上使用了龙门架构型的机械臂完成多个秧苗的同时夹取与放置，在R1后部采用了一组双自由度机械臂配合蓄能弹射机构用于将谷物传递至三区。在R2上使用了一个四自由度机械臂并配合涵道风扇完成对谷物的吸取与放置。
\end{itemize}

团队在2024年ROBOCON全国赛中获“颗粒归仓”赛项\textbf{竞技赛全国三等奖与操作技能挑战赛全国二等奖}。此后以R1机器人为主体的“自主饮料配送机器人”项目获得2024年中国大学生机械工程创新创意大赛\textbf{全国三等奖与长三角赛区一等奖}。

\begin{figure}[H]
    \centering
    \begin{minipage}{0.48\textwidth}
        \centering
        \includegraphics[width=\textwidth]{robocon-2024-2.jpg}
        \caption*{比赛中的R1机器人}
    \end{minipage}
    \hfill
    \begin{minipage}{0.48\textwidth}
        \centering
        \includegraphics[width=\textwidth]{robocon-2024-3.jpg}
        \caption*{比赛中的R2机器人}
    \end{minipage}
\end{figure}

\begin{figure}[H]
    \centering
    \includegraphics[width=0.95\textwidth]{robocon-2024-4.jpg}
    \caption*{2024赛季华东理工大学”无贰“战队队员合影}
\end{figure}

相关链接：
\begin{itemize}
    \item 培训视频：\url{https://space.bilibili.com/691149606}
    \item 演示视频1：\url{https://www.bilibili.com/video/BV1w142187c4/?share_source=copy_web&vd_source=1eb461f6bddf0bc2f2101474c44a8a70}
    \item 演示视频2：\url{https://www.bilibili.com/video/BV11x4y1t7Ji/?share_source=copy_web&vd_source=1eb461f6bddf0bc2f2101474c44a8a70}
    \item 开源图纸：\url{https://pan.baidu.com/s/1ud0ZduSVzKmUnaYjAMiyBw?pwd=z9if}
\end{itemize}
